\section{Conclusion}
This study asks how weakening excluded-instrument relevance changes the finite-sample behavior of preserved high-dimensional IVQR procedures when the documented treatment/endogeneity distribution is held fixed. The question is narrower than whether IVQR can accommodate endogeneity or whether orthogonal scores can support inference under suitable assumptions. Those are questions addressed by the motivating literature. The present contribution is a controlled empirical comparison that traces the relationship among point accuracy, retention of the true candidate, rejection of specified false candidates and localization of the accepted numerical profile. It extends the fixed-relevance simulation of \textcite{chen_huang_tien_2021} without introducing a new estimator or asymptotic theorem.

The experimental construction reduces the excluded signal and adds independent compensating noise under the stated Gaussian assumptions. This preserves the documented joint treatment/outcome-disturbance distribution and the structural effects while changing treatment--instrument relationships. It avoids interpreting a simultaneous change in treatment dispersion as a consequence of relevance alone. The construction does not preserve each realized treatment or outcome across settings, and the relevance parameter is not a universal measure of quantile-specific identification. Its usefulness lies in making a particular comparison explicit and interpretable, rather than in defining a general weak-instrument classification for IVQR applications.

The central empirical result is disagreement among measures that can each appear reassuring in isolation. As relevance weakens, truth acceptance can approach the nominal benchmark while false-candidate rejection becomes limited and the maintained grid retains little localization. The point-error comparisons provide a second warning: lower RMSE at the larger sample size does not imply that truth acceptance moves closer to its benchmark. Strong-relevance distortions and quantile asymmetries remain important counterpoints, showing that weak relevance alone cannot account for all observed patterns. The contribution is the documented configuration of these findings in the controlled experiment, rather than the general observation that weak identification can reduce power.

The benchmark comparisons sharpen that interpretation without identifying its mechanism. Oracle-GMM also deteriorates despite knowing the designated relevant controls, so high-dimensional nuisance estimation cannot be the sole explanation for the loss of performance. DML-IVQR often has smaller point errors than Full-GMM, but the retained exceptions and differences in nuisance dimension, residualization and fitting conventions prevent a universal ranking or an attribution to regularization alone. The comparison concerns complete historical implementations. It does not estimate the effect of switching a single computational component, and it supplies no evidence about methods or DGPs that were not included.

The inferential objects must also remain distinct. A direct-at-truth evaluation describes a candidate decision under the implemented reference cutoff. For a coherently specified randomized test and its inversion, that acceptance event can represent membership of truth in the conceptual inverted set; the identity alone says nothing about attainment of a nominal level. The stored finite-grid accepted set in this experiment is evaluated separately, and fresh pivotal penalties can distinguish direct and stored-profile decisions. Raw false-candidate rejection is not a size-adjusted power comparison. Neither high truth acceptance nor a compact accepted profile establishes theoretically valid confidence-set coverage for the implemented procedures.

Several limitations qualify the theoretical interpretation. The non-finite mixed residual density prevents use of the displayed standard finite-density argument, while the historical code uses a different fitted-normal proxy. Same-sample nuisance estimation, intercept and projection conventions, penalty randomness and finite solver outputs remain features of the objects being compared. No separate experiment isolates their contributions to the rejection patterns. The selected numerical audits document rank and decision sensitivity for specified problems but do not certify every historical decision or a full-experiment error frequency. These qualifications restrict inference claims without invalidating the descriptive content of the frozen numerical records.

The computational domain creates an additional limit. Accepted measure and full-grid frequency describe the investigated grid, not continuous boundaries or unrestricted set geometry. Wider-domain diagnostics reveal how much localization changes over the examined ranges, but do not establish global boundedness or nonidentification. Exact historical stochastic replay is also not established from the parent seed: worker-specific penalty streams and portable installation details are incomplete. These limits explain why evidence inspection, reporting from frozen outputs and numerical replay should be described separately. A transparent computational account is part of interpreting the experiment, not a substitute for an inference argument.

Three concrete research directions follow, without being carried out in this revision. First, a separately specified design with an analytically checked finite mixed density could distinguish the relevance comparison from this particular regularity limitation while preserving clearly stated distributional controls. Second, a component comparison could vary nuisance fitting, residualization, density treatment or cross-fitting under a prespecified randomization scheme, allowing mechanisms that are confounded here to be examined separately. Third, a domain and computation study could predefine wider candidate searches and archive worker states and candidate-level inputs, distinguishing numerical decision stability from inferential calibration. Each would create new evidence and would require its own scientific justification; none retrospectively repairs the present archive.

The immediate scientific implication is therefore one of interpretation. In a simulation of an implemented IVQR procedure, point accuracy, truth acceptance and the ability to exclude alternatives are distinct properties. Reporting them together can expose weak discrimination and poor localization that a single performance measure conceals. Within the studied designs, this controlled comparison provides such evidence while leaving the theoretical validity and causal source of the observed decision patterns unestablished. That combination of substantive findings and explicit limits is the contribution of the article.
